\documentclass[11pt,a4paper]{article}
\usepackage[T1]{fontenc}
\usepackage{lmodern,microtype,amsmath,amssymb,amsthm,mathtools,booktabs,array}
\usepackage[margin=25mm]{geometry}
\usepackage[colorlinks=true,linkcolor=blue,citecolor=blue,urlcolor=blue]{hyperref}
\newtheorem{theorem}{Theorem}
\newtheorem{lemma}[theorem]{Lemma}
\newtheorem{corollary}[theorem]{Corollary}
\newtheorem{proposition}[theorem]{Proposition}
\newcommand{\E}{\mathbb E}
\newcommand{\Prob}{\mathbb P}
\newcommand{\causal}{\mathrm{causal}}
\newcommand{\R}{\mathbb R}
\title{Exact Causal Sources:\\Spectral Realization, Sharp Asymptotics,\\and Stability of Near-Optimal Seeds}
\author{}
\date{4 October 2026}
\begin{document}
\maketitle
\begin{abstract}
For a finite controlled history tree, we construct one exact causal seed
whose information is stochastically between the lower envelope of all
adaptive transcript information laws and that envelope plus an independent
exponential variable of mean $\log_2e$. This yields simultaneous
positive-order Renyi approximation bounds, exact min-entropy, a constant
gap to unconstrained coupling of all adaptive transcripts, and an exact
Wasserstein identity for the entropy excess. For any fixed finite family
of memoryless laws of common entropy $h$, we prove
$C_n^{\causal}=nh+\sqrt{2/\pi}(\sigma_{\max}-\sigma_{\min})\sqrt n+o(\sqrt n)$,
including zero lower variance. Every source with entropy within
$o(\sqrt n)$ of optimal has the same normalized information limit in
$W_1$. The greedy source additionally has convergence of all polynomial
moments and asymptotically optimal fixed-confidence information quantiles.
The square-root analysis uses explicit discrete Bellman barriers.
A separately identified martingale theorem supplies the logarithmic upper
bound when all positive varentropies agree. An irrational crash pair has
an explicit finite-horizon error certificate. The literature map separates
established coupling and nonlinear-CLT results from their compatible-tree
realization and the consequences proved here.
\end{abstract}
\tableofcontents
\newpage

\section{Position in the literature and the scope of the claims}\label{sec:literature}
This note concerns an exact, finite-horizon, history-indexed controlled
interface. One random seed must implement the conditional output law at
every history and for every requested action. The distinction between this
problem and an unconstrained coupling of transcript laws is essential.
All logarithms in information quantities are base two; $\ln$ is natural.

\paragraph{Ordinary minimum-entropy coupling.}
Kocaoglu, Dimakis, Vishwanath and Hassibi \cite{KDVH17} study greedy
minimum-entropy coupling in connection with entropic causal inference.
Li \cite{Li21} develops approximation methods for multiple marginals.
Compton \cite{C22} proves the $\log_2e$ additive guarantee for the ordinary
greedy coupling algorithm. Compton et al.\ \cite{CKQGK23} improve the
Shannon guarantees to $\log_2(e)/e$ for two marginals and
$(1+\log_2e)/2$ for arbitrarily many marginals, using their profile method.
These results concern fixed marginal distributions. Their algorithms do
not, by their statements alone, provide a compatible decoder at shared
history-action nodes of a controlled tree. The tree construction below
supplies that additional compatibility directly. Its $\log_2e$ bound
is not asserted to be an improvement of those ordinary-coupling constants.

\paragraph{Information-spectrum converses and Renyi entropy.}
Shkel and Yadav \cite{SY23}, Theorem 1 and Section III-B, give
information-spectrum lower bounds for functional representations and
minimum Shannon/Renyi entropy couplings. The lower envelope used here
is an application of that spectrum viewpoint to all deterministic
adaptive-policy transcript laws. The converse is reproved below.
Ma, Wang and Wu \cite{MWW25} establish parameter-dependent greedy
approximation guarantees for ordinary two-marginal minimum-Renyi
coupling; their error tends to zero as the order tends to infinity.
Consequently neither ``Renyi coupling'' nor ``a vanishing high-order
approximation gap'' should be advertised as new in itself.
Our bound is derived from one compatible tree seed, simultaneously for
all orders, with an explicit elementary error function.
The recent extreme-point work of Ma et al.\ \cite{MWWC25} addresses exact
ordinary coupling through a different finite-dimensional route.

\paragraph{Controlled central limits.}
Peng's nonlinear central limit theory \cite{Peng07} and subsequent
quantitative work \cite{FPSS18} provide the established setting of
uncertain conditional variance and nonlinear heat equations.
Peng and Zhou \cite{PZ19}, Theorem 2, give an explicit one-sided tail
capacity. With lower and upper standard deviations $b,a$, its complement
is precisely the CDF $F_{a,b}$ used below. We do not claim that this
profile or the underlying extremal-variance principle is new.
We give a direct discrete comparison for arbitrary finite increment
laws and then transfer it to the information of an exact causal seed.
The CDF envelope defines an ordinary probability law, but this must not
be confused with representing the entire nonlinear expectation by a
single probability measure. Different thresholds can have different
optimizing policies; when $a>b$, the envelope has positive mean even
though every individual centered transcript is a martingale sum.

\paragraph{Equal conditional variances.}
The only substantive external estimate used in the proofs is the
nonuniform martingale Berry--Esseen inequality of Fan, Grama and Liu
\cite{FGL}, Theorem 2.1. We state exactly the specialization needed and
check its hypotheses. Its unnamed absolute constant is retained as
$C_{\rm FGL}$; we do not assign it an unproved numerical value.
All other arguments below are given directly.

\paragraph{What this manuscript establishes, and what it does not.}
The central finite-tree result is a stochastic sandwich for the information
of one compatible greedy seed. Its consequences include simultaneous
Renyi bounds, exact min-entropy, a constant gap to the ordinary coupling
of \emph{all adaptive} transcript laws, and a Wasserstein stability identity.
The latter forces a common information limit for every first-order
near-optimal seed on the square-root scale. We also prove convergence
of all polynomial moments for the greedy construction and explain why
that stronger conclusion fails for arbitrary near-optimal seeds.
The literature map checks the primary sources cited above, including
2025 work; it is not an exhaustive priority certificate. We make no
claim to have discovered a new $G$-normal law or to have settled the
sharp logarithmic regime.


\section{Model and the sharp square-root statement}\label{sec:model}
Fix a finite nonempty action family $\mathcal A=\{1,\ldots,m\}$.
Action $i$ requests a finite law $P_i$; zero-probability symbols are
removed. Assume in the asymptotic sections that
\[
H(P_i)=h\quad(i\in\mathcal A),\qquad
\sigma_i^2=\operatorname{Var}_{P_i}(-\log_2P_i(Y)),\qquad
a=\max_i\sigma_i,\quad b=\min_i\sigma_i.
\]
A seed $S$ is sampled once, independently of external policy randomness.
At round $t$, the output is a deterministic function of $S$, the complete
past action-output history, and the currently requested action.
Under every adaptive policy its conditional law must be the requested
$P_i$. Let $C_n^{\causal}$ be the infimum Shannon entropy of finite seeds
with this property. Seeds may depend on $n$. Responses to the same action
at distinct histories need not agree. Thus this is a history-indexed
causal model, not the stronger synchronous model in which a time-action
coordinate must be shared across different histories.

\begin{theorem}\label{thm:main}
For every fixed finite family as above,
\begin{equation}\label{eq:main}
\boxed{C_n^{\causal}
=nh+\sqrt{\frac2\pi}(a-b)\sqrt n+o(\sqrt n).}
\end{equation}
If $a=b=0$, then $C_n^{\causal}=nh$. If $a=b>0$, then
\begin{equation}\label{eq:equalvar}
0\le C_n^{\causal}-nh=O(\log(n+1)).
\end{equation}
\end{theorem}
Write $\nu_i$ for the law of $D_i=-\log_2P_i(Y)-h$, so $\E D_i=0$.
Define
\begin{equation}\label{eq:T}
(Tf)(x)=\min_{i\in\mathcal A}\E f(x-D_i),\qquad
G_0(x)=\mathbf1_{\{x\ge0\}},\qquad G_n=T^nG_0.
\end{equation}
Each $G_n$ is a finite-support CDF; let $Z_n\sim G_n$, $e_n=\E Z_n$.
Here $Z_n$ is centered relative to the entropy baseline; the uncentered
information envelope below has law $nh+Z_n$.
For a deterministic policy $\pi$, let $M_{\pi,n}$ be transcript
information minus $nh$. Backward induction gives
\begin{equation}\label{eq:policy}
G_n(x)=\min_\pi\Prob(M_{\pi,n}\le x).
\end{equation}
After the first action and observed output, a separate continuation may
be chosen for each branch and remaining threshold. This gives precisely
\eqref{eq:T}. Extra observed labels with the same increment do not change
the optimal value: each continuation has the same available kernels and
threshold. A randomized policy gives a mixture of deterministic-policy
values after conditioning on its private randomness, so cannot lower
the minimum further. Deterministic actions contribute no additional
information beyond their preceding history.

The proof will establish
\begin{equation}\label{eq:twoingredients}
0\le C_n^{\causal}-nh-e_n\le\log_2e,
\qquad e_n/\sqrt n\longrightarrow\sqrt{2/\pi}(a-b).
\end{equation}

\section{Global realization with a constant entropy gap}
We prove a finite-tree statement, also valid for history-dependent
kernels and any finite number of actions. A history is $v$, and its
child after action $A$ and output $y$ is $vAy$. Set
\[
p(\varnothing)=1,\qquad p(vAy)=p(v)K_v(y\mid A).
\]
These cylinder weights satisfy
$\sum_y p(vAy)=p(v)$ for each action separately. For a policy $\pi$,
let $\mathcal L_\pi$ denote its terminal histories. Their weights sum
to one and form its transcript law. Define $I_\pi=-\log_2p(T_\pi)$ and
\[
F_*(t)=\min_\pi\Prob(I_\pi\le t),\qquad Z\sim F_*,\quad L=\E Z.
\]

There are finitely many deterministic response maps on this finite tree.
Any seed can be coarsened by merging atoms with the same response map,
without changing its outputs or increasing a positive-order Renyi entropy.
The feasible weights on these maps form a nonempty compact polytope:
nonemptiness follows, for example, by independently pre-sampling a response
at every history-action node. Positive-order Renyi entropies, Shannon
entropy, and min-entropy are continuous on this simplex, so their minima
are attained. We continue to use infimum notation when convenient.

\begin{lemma}[Finite global realization]\label{lem:realization}
There is a finite exact causal seed $S_G$, supported on at most the
number of terminal histories in the full tree, such that
\begin{equation}\label{eq:gap}
L\le C\le H(S_G)\le L+\log_2e.
\end{equation}
\end{lemma}
\begin{proof}
A residual flow $r$ is nonnegative and satisfies
$\sum_y r(vAy)=r(v)$ for every action. A deterministic response
strategy $d$ chooses an output at each history-action pair. Its
compatibility indicator $\chi_d(v)$ equals one when all its outputs
on the path to $v$ agree with $v$. It is a flow of root mass one.

For any residual flow compute from the leaves upward
\[
B_r(v)=r(v)\quad(v\text{ terminal}),\qquad
B_r(v)=\min_A\max_y B_r(vAy)\quad(v\text{ nonterminal}).
\]
There are two elementary consequences.
First, the largest extractable strategy weight is $B_r(\varnothing)$:
for each possible action the strategy selects a child of maximal
bottleneck, and this is done recursively. Nonnegativity on terminal
histories implies nonnegativity at all histories by flow conservation.
Second, a policy choosing a minimizing action at every history has
\[
r(v)\le B_r(\varnothing)\qquad(v\in\mathcal L_{\pi_r}).
\]
Indeed, along every output of that policy, bottleneck values never
increase. This policy also proves that no larger strategy weight is
possible. If the root mass is positive, every action has a positive
child, and induction gives a positive root bottleneck.

Starting from $r=p$, repeatedly extract a maximizing strategy of weight
$w=B_r(\varnothing)$ and replace $r$ by $r-w\chi_d$. At least one
positive terminal residual becomes zero at each extraction; otherwise
the same strategy would permit a larger weight. Residuals never
increase, so the algorithm terminates after finitely many steps. Its
weights $w_j$ are nonincreasing, sum to one, and satisfy
\begin{equation}\label{eq:decomp}
p(v)=\sum_jw_j\chi_{d_j}(v).
\end{equation}
Choosing strategy $d_j$ with probability $w_j$ gives all the prescribed
history weights simultaneously. Ratios of child to parent weights
give the correct conditional laws at every positive-probability history.
Thus this is one exact causal seed for all policies.

Fix $0<\delta\le1$ and stop just after all weights larger than $\delta$
have been extracted. The remaining root mass is
$r(\varnothing)=\sum_{j:w_j\le\delta}w_j$. Its bottleneck is at most
$\delta$. The minimizing policy above therefore gives
\begin{align}
\sum_{j:w_j\le\delta}w_j
&\le\max_\pi\sum_{v\in\mathcal L_\pi}\min\{p(v),\delta\}\notag\\
&=\max_\pi\E\min\{1,\delta2^{I_\pi}\}
\le\E\min\{1,\delta2^Z\}.\label{eq:softtail}
\end{align}
The last inequality uses stochastic dominance: $F_*\le F_{I_\pi}$.
The zero residual case satisfies the same inequality trivially.
For $J=-\log_2w_{S_G}$, substitute $\delta=2^{-t}$ to obtain
\[
\Prob(J\ge t)\le\E\min\{1,2^{Z-t}\}
=\Prob(Z+E\ge t),
\]
where $E$ is independent exponential of rate $\ln2$. Integrating these
nonnegative-variable tails yields $H(S_G)\le L+1/\ln2$.

For the converse, under any deterministic policy each seed atom is
contained in its transcript atom. Its surprisal therefore dominates
the transcript surprisal under the induced coupling. Hence
$F_S\le F_{I_\pi}$ for all $\pi$, and thus $F_S\le F_*$. Integration
gives $H(S)\ge L$. This proves \eqref{eq:gap}.
\end{proof}

In the equal-entropy memoryless model, $F_*(t)=G_n(t-nh)$ by
\eqref{eq:policy}. Lemma~\ref{lem:realization} proves the first
inequality in \eqref{eq:twoingredients}. In particular, realizability
introduces no error proportional to the number of rounds or blocks.


\section{Finite-tree consequences: all Renyi orders and stability}
For an exact seed $S$, write $J_S=-\log_2\Prob(S)$. The envelope variable
in Lemma~\ref{lem:realization} is denoted by $Z$ and need not itself be
the information of any discrete source. Its law is nevertheless a
well-defined finite-support probability distribution.

\begin{theorem}[Stochastic realization]\label{thm:sandwich}
For every finite controlled history tree and its greedy seed $S_G$,
\begin{equation}\label{eq:sandwich}
Z\preceq_{\rm st}J_{S_G}\preceq_{\rm st}Z+E,
\qquad E\sim\operatorname{Exp}(\ln2),\quad E\perp Z.
\end{equation}
Every exact seed satisfies $Z\preceq_{\rm st}J_S$.
The notation $X\preceq_{\rm st}Y$ means $F_X(t)\ge F_Y(t)$ for all $t$.
\end{theorem}
\begin{proof}
The atom-containment argument in Lemma~\ref{lem:realization} proves the
lower stochastic bound for every seed. Its inequality
\eqref{eq:softtail}, with $\delta=2^{-t}$, gives
$\Prob(J_{S_G}\ge t)\le\E\min(1,2^{Z-t})=\Prob(Z+E\ge t)$
for $t\ge0$. All variables are nonnegative, so the comparison also holds
for $t<0$. Taking one-sided limits supplies the CDF formulation,
including at atoms.
\end{proof}
This is an ordering of distributions. It does not say that the algorithm
constructs $J_{S_G}=Z+E$, or that a seed's excess is an independent
exponential variable on its operational probability space.

For $0<\alpha<\infty$, $\alpha\ne1$, define
\[
H_\alpha(S)=\frac{1}{1-\alpha}\log_2\sum_s\Prob(S=s)^\alpha,
\qquad R_\alpha(Z)=\frac{1}{1-\alpha}\log_2\E2^{(1-\alpha)Z}.
\]
Set $H_1=H$, $R_1(Z)=\E Z$, and let $C_\alpha$ be the infimum of
$H_\alpha$ over exact causal seeds for this tree. Write
\[
g_\alpha=\frac{\log_2\alpha}{\alpha-1}\quad(\alpha\ne1),
\qquad g_1=\log_2e.
\]

\begin{theorem}[A single seed for every positive Renyi order]\label{thm:renyi}
Simultaneously for all $0<\alpha<\infty$,
\begin{equation}\label{eq:renyi}
R_\alpha(Z)\le C_\alpha\le H_\alpha(S_G)
\le R_\alpha(Z)+g_\alpha.
\end{equation}
The construction of $S_G$ is independent of $\alpha$.
\end{theorem}
\begin{proof}
Put $M_\alpha=\E2^{(1-\alpha)Z}$. Independence and direct integration give
\[
\E2^{(1-\alpha)(Z+E)}=M_\alpha/\alpha,
\qquad
\E2^{(1-\alpha)E}=
\int_0^\infty (\ln2)e^{-\alpha(\ln2)t}\,dt=1/\alpha.
\]
For $0<\alpha<1$, the function $2^{(1-\alpha)x}$ is increasing, so
\eqref{eq:sandwich} yields
$M_\alpha\le\sum_s\Prob(S_G=s)^\alpha\le M_\alpha/\alpha$.
Take logarithms and divide by $1-\alpha>0$.
For $\alpha>1$, the moment inequalities reverse; division by
$1-\alpha<0$ reverses them again, giving the same bound.
For arbitrary $S$, the lower stochastic bound proves the lower entropy
bound by the same sign check. The case $\alpha=1$ is
Lemma~\ref{lem:realization}, or follows by taking limits.
\end{proof}

\begin{theorem}[Exact min-entropy]\label{thm:minentropy}
Let $B_p$ be the bottleneck recursion in Lemma~\ref{lem:realization}.
Then
\[
C_\infty=H_\infty(S_G)=-\log_2 B_p(\varnothing).
\]
For memoryless laws $P_1,\ldots,P_m$ repeated for $n$ rounds, without
any equal-entropy assumption,
\begin{equation}\label{eq:minentropy}
C_{n,\infty}=n\max_iH_\infty(P_i).
\end{equation}
\end{theorem}
\begin{proof}
An atom of any exact seed determines a pure response strategy. Its
weight cannot exceed that strategy's extractable weight, and hence
cannot exceed $B_p(\varnothing)$. The first greedy atom has exactly
that weight, and all later weights are no larger. This proves the first
claim. Put $\beta=\min_i\max_yP_i(y)$. Induction on remaining depth
shows $B_p(v)=p(v)\beta^{n-|v|}$ in the memoryless model. Taking minus
base-two logarithms proves \eqref{eq:minentropy}.
\end{proof}
No finite bound at order $\alpha=0$ follows by sending $\alpha\downarrow0$
in \eqref{eq:renyi}: $g_\alpha$ diverges. The number of atoms is instead
bounded by the number of full-tree terminal histories.

\begin{theorem}[Cost of compatibility with all adaptive transcripts]\label{thm:compatibility}
Let $M_\alpha$ denote the minimum Renyi entropy of an \emph{ordinary}
joint coupling of the transcript laws for all deterministic adaptive
policies in the finite tree. Then
\[
0\le C_\alpha-M_\alpha\le g_\alpha\quad(0<\alpha<\infty),
\qquad C_\infty=M_\infty.
\]
\end{theorem}
\begin{proof}
An exact causal seed induces a joint vector of transcripts by running
each deterministic policy on that same seed. This vector is a
deterministic function of the seed. Merging probability atoms cannot
increase any positive-order Renyi entropy, so $M_\alpha\le C_\alpha$.
Conversely, the information of any ordinary joint coupling dominates
the information of each of its coordinates. It therefore dominates
$Z$ stochastically and has entropy at least $R_\alpha(Z)$.
Theorem~\ref{thm:renyi} now gives
$C_\alpha\le R_\alpha(Z)+g_\alpha\le M_\alpha+g_\alpha$.
For the infinite-order endpoint, the largest atom in any such coupling
is at most $\min_\pi\max_{v\in\mathcal L_\pi}p(v)$.
The minimizing bottleneck policy and the maximizing response strategy
prove that this min--max is $B_p(\varnothing)$; apply
Theorem~\ref{thm:minentropy}.
\end{proof}
All adaptive policies are required here. Replacing them by constant
policies, or by one-time output marginals, changes the lower bound.
This is a finite-horizon comparison, not a claim that a single seed
works for every horizon.

\begin{corollary}[Tail and quantile guarantees]\label{cor:quantiles}
For all real $t$ and $s\ge0$,
\[
\Prob(J_{S_G}>t+s)\le\Prob(Z>t)+2^{-s}.
\]
Writing $q_p(X)=\inf\{x:F_X(x)\ge p\}$, for $0<\delta<\epsilon<1$,
\[
q_{1-\epsilon}(Z)\le q_{1-\epsilon}(J_{S_G})
\le q_{1-\epsilon+\delta}(Z)+\log_2(1/\delta).
\]
\end{corollary}
\begin{proof}
Use \eqref{eq:sandwich} and
$\{Z+E>t+s\}\subset\{Z>t\}\cup\{E>s\}$.
Substitute $t=q_{1-\epsilon+\delta}(Z)$ and
$s=\log_2(1/\delta)$ for the upper quantile bound. Stochastic dominance
gives the lower one. Finite support ensures the defining CDF levels
are attained.
\end{proof}

\begin{lemma}[Exact transport identity]\label{lem:transport}
For every finite-entropy exact seed $S$,
\begin{equation}\label{eq:transport}
W_1(\mathcal L(J_S),\mathcal L(Z))=H(S)-\E Z.
\end{equation}
In particular, if $H(S)\le C_1+\Delta$, then this distance is at most
$\Delta+\log_2e$.
\end{lemma}
\begin{proof}
For any coupling, $\E|J_S-Z|\ge|\E J_S-\E Z|$.
For a uniform $U\in(0,1)$, stochastic dominance gives
$q_U(J_S)\ge q_U(Z)$. This quantile coupling attains the lower bound:
its expected absolute difference is $\E J_S-\E Z$.
The last assertion follows from Lemma~\ref{lem:realization}.
\end{proof}

\begin{theorem}[Transfer of any diverging-scale envelope limit]\label{thm:transfer}
Consider a sequence of finite trees, envelope variables $Z_n$, centers
$m_n$, and scales $r_n\to\infty$. Suppose
$\mathcal L((Z_n-m_n)/r_n)\to\mu$ in $W_1$.
Every exact seed sequence satisfying $H(S_n)-C_{n,1}=o(r_n)$ then obeys
\[
\mathcal L((J_{S_n}-m_n)/r_n)\longrightarrow\mu\quad\hbox{in }W_1.
\]
\end{theorem}
\begin{proof}
Translation invariance, scaling, and Lemma~\ref{lem:transport} show that
the distance between the two normalized information laws is bounded by
$(H(S_n)-C_{n,1}+\log_2e)/r_n\to0$. Apply the triangle inequality.
\end{proof}

\section{An explicit finite-horizon Bellman comparison}
Assume for now $a\ge b>0$. Let $\Phi$ be the standard normal CDF and set
\begin{equation}\label{eq:F}
F(x)=\begin{cases}
\displaystyle\frac{2b}{a+b}\Phi(x/b),&x\le0,\\[5pt]
\displaystyle1-\frac{2a}{a+b}\bigl(1-\Phi(x/a)\bigr),&x>0.
\end{cases}
\end{equation}
This is the complement of the one-sided tail capacity in
Peng--Zhou \cite{PZ19}, Theorem 2. We verify the profile directly. Define
\[
K=\frac{\sqrt{2/\pi}}{a+b},\quad
m_3=\max_A\E|D_A|^3,\quad
\mathcal C=K\left(\frac{m_3}{3b^2}+a\right),\quad
u_t(x)=F(x/\sqrt t)\quad(t>0).
\]

\begin{lemma}[Explicit comparison]\label{lem:comparison}
For every $\tau\ge1$, $R\ge0$, and integer $n\ge0$, put
\[
\eta=\max\{F(-R/\sqrt\tau),\,1-F(R/\sqrt\tau)\}.
\]
Then, for every $x\in\R$,
\begin{equation}\label{eq:barrier}
\boxed{\begin{aligned}
F\!\left(\frac{x-R}{\sqrt{n+\tau}}\right)
-\eta-\frac{3\mathcal C}{\sqrt\tau}
&\le G_n(x),\\
G_n(x)&\le
F\!\left(\frac{x+R}{\sqrt{n+\tau}}\right)
+\eta+\frac{3\mathcal C}{\sqrt\tau}.
\end{aligned}}
\end{equation}
\end{lemma}
\begin{proof}
For $x$ of either sign let $c=b$ on the negative half-line and $c=a$
on the positive half-line. Direct differentiation gives
\[
F'(x)=Ke^{-x^2/(2c^2)},\qquad
F''(x)=-\frac{Kx}{c^2}e^{-x^2/(2c^2)}.
\]
The first two derivatives extend continuously across zero; in particular
$F''(0)=0$. The one-sided third derivatives satisfy
\[
|F'''(x)|=\frac K{c^2}
\left|\frac{x^2}{c^2}-1\right|e^{-x^2/(2c^2)}
\le\frac{2K}{b^2}.
\]
Consequently $F''$ is globally Lipschitz with that constant. Taylor's
theorem therefore applies even when an increment crosses the interface.

The function $u_t$ satisfies the following identity everywhere,
including $x=0$:
\begin{equation}\label{eq:pde}
\partial_tu_t(x)
=\frac12\min\{a^2\partial_{xx}u_t(x),b^2\partial_{xx}u_t(x)\}.
\end{equation}
The minimum over $i$ in the quadratic term uses only the extreme
variances. On $x<0$ the second derivative is positive and the minimum selects
$b^2$; on $x>0$ it is negative and selects $a^2$. On each side both
expressions equal $-x\partial_xu_t(x)/(2t)$.

Because $\E D_A=0$, the spatial Taylor remainder gives
\[
\left|Tu_t-u_t-\partial_tu_t\right|
\le\frac{Km_3}{3b^2}\,t^{-3/2}
\qquad\text{uniformly in }x.
\]
For clarity, the minimum of the finitely many quadratic Taylor terms is exactly
\eqref{eq:pde}; taking a minimum cannot increase the maximal absolute
remainder. Also, with $z=x/\sqrt t$,
\[
\partial_{tt}u_t(x)
=\frac K{4t^2}(3z-z^3/c^2)e^{-z^2/(2c^2)},
\qquad \|\partial_{tt}u_t\|_\infty\le\frac{2Ka}{t^2}.
\]
The last bound follows from
$\sup_{s\ge0}se^{-s^2/2}\le1$ and
$\sup_{s\ge0}s^3e^{-s^2/2}\le2$.
The time Taylor remainder is thus at most $Ka t^{-2}$. For $t\ge1$,
\begin{equation}\label{eq:local}
\|Tu_t-u_{t+1}\|_\infty\le\mathcal C t^{-3/2}.
\end{equation}

The operator $T$ is monotone, commutes with spatial translations and
addition of constants, and is nonexpansive in the sup norm. Telescoping
\eqref{eq:local} therefore gives
\begin{equation}\label{eq:telescoping}
\|T^nu_\tau-u_{n+\tau}\|_\infty
\le\mathcal C\sum_{j=0}^{n-1}(\tau+j)^{-3/2}
\le\frac{3\mathcal C}{\sqrt\tau}.
\end{equation}
Finally, the definition of $\eta$ implies the pointwise initial bounds
\[
u_\tau(x-R)-\eta\le\mathbf1_{\{x\ge0\}}
\le u_\tau(x+R)+\eta.
\]
Apply $T^n$, use its three invariances, and then
\eqref{eq:telescoping}. This proves \eqref{eq:barrier}.
\end{proof}

\begin{corollary}[CDF limit when $b>0$]\label{cor:cdf}
For each $x\in\R$, $G_n(x\sqrt n)\to F(x)$.
\end{corollary}
\begin{proof}
Fix $\tau,R$ and substitute $x\sqrt n$ in \eqref{eq:barrier}. Both
shifted arguments tend to $x$, so the limsup and liminf differ from
$F(x)$ by at most $\eta+3\mathcal C/\sqrt\tau$.
First choose $\tau$ large and then $R$ large. The error tends to zero.
For explicit choices at tolerance $0<\epsilon<1$, take
\[
\tau\ge\max\{1,(6\mathcal C/\epsilon)^2\},\qquad
R=a\sqrt{2\tau\ln(4/\epsilon)}.
\]
The Gaussian Chernoff bound gives $\eta\le\epsilon/2$, and the other
term is at most $\epsilon/2$. These parameters are fixed before sending
$n$ to infinity, so no uniform convergence over a moving policy grid
is being assumed.
\end{proof}


\section{Zero variance and convergence of the means}
\begin{lemma}[The degenerate endpoint]\label{lem:zero}
If $a>0$ and $b=0$, then
\[
G_n(x\sqrt n)\longrightarrow F_{a,0}(x)
=\begin{cases}0,&x\le0,\\2\Phi(x/a)-1,&x>0.\end{cases}
\]
\end{lemma}
\begin{proof}
This step concerns the centered increment control problem only.
For $\varepsilon>0$ add an independent Rademacher increment
$\varepsilon\xi_t$ at each time and for every action. Retain both the
original output and $\xi_t$ as observed labels. The new centered
increments have extreme variances $a^2+\varepsilon^2$ and $\varepsilon^2$;
write $G_n^\varepsilon$ for their Bellman CDF.
Extra labels do not change the Bellman recursion: conditional on the
history the next increment law depends only on the action, and the
remaining terminal threshold depends only on the current sum.

The perturbed and unperturbed sums under the same action choices differ
by $\varepsilon\sum_{t=1}^n\xi_t$. For every $d>0$, Chebyshev gives
\[
\Prob\left(\left|\varepsilon\sum_{t=1}^n\xi_t\right|>d\sqrt n\right)
\le\varepsilon^2/d^2,
\]
uniformly over policies. Every original policy is a perturbed policy
that ignores the extra coins. Conversely, an original controller can
emulate a perturbed policy by drawing those coins as private randomness.
The optimal original value is no larger than that randomized policy's
value. These two observations give
\begin{equation}\label{eq:noise}
G_n^\varepsilon((x-d)\sqrt n)-\varepsilon^2/d^2
\le G_n(x\sqrt n)
\le G_n^\varepsilon((x+d)\sqrt n)+\varepsilon^2/d^2.
\end{equation}
Apply Corollary~\ref{cor:cdf} with
$a_\varepsilon=\sqrt{a^2+\varepsilon^2}$ and $b_\varepsilon=\varepsilon$.
Then let $\varepsilon\downarrow0$ with $d=\sqrt\varepsilon$.
The explicit CDFs \eqref{eq:F} converge to the continuous CDF $F_{a,0}$,
also at the moving arguments $x\pm\sqrt\varepsilon$. At $x=0$, both
limits are zero directly from the two formulas. Equation~\eqref{eq:noise}
proves the claim.
\end{proof}

\begin{lemma}[Mean limit]\label{lem:mean}
For $a\ge b\ge0$,
\begin{equation}\label{eq:mean}
\lim_{n\to\infty}\frac{e_n}{\sqrt n}
=\sqrt{\frac2\pi}\,(a-b).
\end{equation}
\end{lemma}
\begin{proof}
Let $B=\max_{A,y}|-\log_2P_A(y)-h|$. If $B=0$ the claim is immediate;
assume $B>0$. For any policy the centered
increments have conditional mean zero and absolute value at most $B$.
Convexity bounds their conditional exponential moment by
$\cosh(\lambda B)\le e^{\lambda^2B^2/2}$. Iteration and exponential
Markov therefore give, uniformly over all policies,
\[
\Prob(M_{\pi,n}>t\sqrt n),\quad
\Prob(M_{\pi,n}\le-t\sqrt n)
\le e^{-t^2/(2B^2)}\qquad(t>0).
\]
By \eqref{eq:policy}, both tails of $Z_n/\sqrt n$ have these same
bounds. Hence the signed tail formula
\[
\frac{e_n}{\sqrt n}
=\int_0^\infty[1-G_n(t\sqrt n)]\,dt
-\int_0^\infty G_n(-t\sqrt n)\,dt
\]
admits dominated convergence, using Corollary~\ref{cor:cdf} or
Lemma~\ref{lem:zero}. When $b>0$, its limit is
\begin{align*}
\int_0^\infty[1-F(t)]\,dt-\int_0^\infty F(-t)\,dt
&=\frac{2a^2}{(a+b)\sqrt{2\pi}}
-\frac{2b^2}{(a+b)\sqrt{2\pi}}\\
&=\sqrt{2/\pi}\,(a-b).
\end{align*}
Here $\int_0^\infty[1-\Phi(t/c)]\,dt=c/\sqrt{2\pi}$.
For $b=0<a$ the negative integral vanishes and the positive one is
$a\sqrt{2/\pi}$. If $a=b=0$, all increment laws are point masses
at zero and $G_n=G_0$, so the assertion is immediate.
\end{proof}

\begin{proof}[Proof of Theorem~\ref{thm:main}]
Lemma~\ref{lem:realization} and Lemma~\ref{lem:mean} give
\[
\frac{C_n^{\causal}-nh}{\sqrt n}
=\frac{e_n}{\sqrt n}+O(n^{-1/2})
\longrightarrow\sqrt{2/\pi}\,(a-b).
\]
If all varentropies vanish, all laws are uniform. Equal entropy means
their supports have the same size $k$. A single uniform seed on
$\{1,\ldots,k\}^n$, with the appropriate output relabeling for each
requested action, has entropy $n\log_2k=nh$. The all-$P$ transcript
gives the reverse inequality.
The stronger bound \eqref{eq:equalvar} is proved in Section~\ref{sec:fgl}.
\end{proof}


\section{Rigidity of near-optimal sources and higher moments}\label{sec:rigidity}
Let $W_{a,b}$ have CDF \eqref{eq:F} when $b>0$, the half-normal
CDF of Lemma~\ref{lem:zero} when $b=0<a$, and the point mass at zero
when $a=b=0$. In this section $Z_n$ is the \emph{centered} Bellman
variable defined in \eqref{eq:T}; the uncentered finite-tree envelope
of Theorem~\ref{thm:sandwich} is $nh+Z_n$.

\begin{lemma}[Envelope convergence in transport]\label{lem:envw1}
The laws of $Z_n/\sqrt n$ converge to $\mathcal L(W_{a,b})$ in $W_1$.
They also converge in every fixed absolute moment.
\end{lemma}
\begin{proof}
The preceding CDF limit gives weak convergence. Lemma~\ref{lem:mean}
bounds both tails of $Z_n/\sqrt n$ by $e^{-t^2/(2B^2)}$, uniformly in $n$.
For any $r>0$, the formula
$\E|X|^r=r\int_0^\infty t^{r-1}\Prob(|X|>t)\,dt$
and dominated convergence give absolute-moment convergence.
For completeness, the one-dimensional quantile coupling identity is
\[
W_1(\mathcal L(X),\mathcal L(Y))
=\int_{\R}|F_X(t)-F_Y(t)|\,dt.
\]
It follows by integrating the lengths of the intervals between the two
quantiles over their common uniform parameter (Tonelli's theorem).
The Gaussian tail majorant, together with pointwise CDF convergence,
therefore proves the $W_1$ statement directly. The deterministic case
$B=0$ is immediate.
\end{proof}

\begin{theorem}[Information-spectrum rigidity]\label{thm:rigidity}
Let $S_n$ be any sequence of exact causal seeds for the fixed finite
family of equal-entropy laws in Theorem~\ref{thm:main}. If
\[
H(S_n)-C_n^{\causal}=o(\sqrt n),
\]
then, writing $J_n=-\log_2\Prob(S_n)$,
\begin{equation}\label{eq:rigidity}
\mathcal L\left(\frac{J_n-nh}{\sqrt n}\right)
\longrightarrow\mathcal L(W_{a,b})\quad\hbox{in }W_1.
\end{equation}
In particular this holds for entropy minimizers and for the greedy seeds.
\end{theorem}
\begin{proof}
Apply Theorem~\ref{thm:transfer} to the uncentered envelope $nh+Z_n$,
center $nh$, and scale $\sqrt n$, using Lemma~\ref{lem:envw1}.
Equivalently, the exact finite-horizon distance is
\[
W_1\left(\mathcal L\left(\frac{J_n-nh}{\sqrt n}\right),
          \mathcal L\left(\frac{Z_n}{\sqrt n}\right)\right)
=\frac{H(S_n)-nh-e_n}{\sqrt n}
\le\frac{H(S_n)-C_n^{\causal}+\log_2e}{\sqrt n}.
\]
The triangle inequality finishes the proof.
\end{proof}
For $a=b=\sigma>0$, $W_{a,b}$ is exactly $N(0,\sigma^2)$.
For $b=0<a$, it is the law of $a|N(0,1)|$. For $a>b>0$, its density
is continuous at zero and consists of Gaussian pieces with different
scales, as obtained by differentiating \eqref{eq:F}.
The assertion concerns the information law, not uniqueness of the seed
alphabet or of its deterministic decoders.

\begin{theorem}[All polynomial moments for the greedy source]\label{thm:moments}
For the greedy information $J_n^G$ and every $r>0$,
\[
\E\left|\frac{J_n^G-nh}{\sqrt n}\right|^r\longrightarrow\E|W_{a,b}|^r.
\]
When $a+b>0$,
\begin{equation}\label{eq:absmoment}
\E|W_{a,b}|^r
=\frac{2^{r/2}\Gamma((r+1)/2)}{\sqrt\pi(a+b)}
\bigl(a^{r+1}+b^{r+1}\bigr).
\end{equation}
In particular,
\begin{equation}\label{eq:variance}
\frac{\operatorname{Var}(J_n^G)}n
\longrightarrow a^2-ab+b^2-\frac2\pi(a-b)^2.
\end{equation}
The right side is zero when $a=b=0$.
\end{theorem}
\begin{proof}
Write $Y_n=(J_n^G-nh)/\sqrt n$. The stochastic sandwich implies, for $t>0$,
\begin{align*}
\Prob(Y_n<-t)&\le e^{-t^2/(2B^2)},\\
\Prob(Y_n>t)&\le\Prob(Z_n/\sqrt n>t/2)+\Prob(E>t\sqrt n/2)\\
&\le e^{-t^2/(8B^2)}+e^{-(\ln2)t/2},\qquad n\ge1.
\end{align*}
These bounds have all polynomial moments and are uniform in $n$.
Combine them with Theorem~\ref{thm:rigidity} and the tail integral formula.
For the moment formula, integrate the density on each half-line and use
\[
\int_0^\infty x^r e^{-x^2/(2c^2)}\,dx
=2^{(r-1)/2}c^{r+1}\Gamma((r+1)/2).
\]
The same formula holds at $b=0$ by the half-normal density.
The mean is $\sqrt{2/\pi}(a-b)$ and the second moment is
$(a^3+b^3)/(a+b)=a^2-ab+b^2$. Subtract its squared mean.
If $B=0$, the exact uniform construction in Theorem~\ref{thm:main}
is also the greedy construction: all extracted weights are $2^{-nh}$.
\end{proof}

\begin{proposition}[Near-optimal entropy does not control variance]\label{prop:rare}
Even within a fixed finite equal-entropy model, there exist exact seeds
$S_n'$ with $H(S_n')-C_n^{\causal}=O(1)$ but
$\operatorname{Var}(-\log_2\Prob(S_n'))/n\to\infty$.
\end{proposition}
\begin{proof}
Append to a greedy seed an independent random variable $R_n$, ignored
by the decoder. For $n\ge2$, let $p_n=n^{-2}$, $M_n=2^{n^2}$, and give
$R_n$ one atom of probability $1-p_n$ and $M_n$ other atoms, each of
probability $p_n/M_n$. Then
\[
H(R_n)=h_2(p_n)+p_n\log_2M_n=h_2(n^{-2})+1=O(1).
\]
Its information has two possible values, and their difference is
$n^2+2\log_2n+\log_2(1-n^{-2})$. Thus
\[
\operatorname{Var}(-\log_2\Prob(R_n))
=p_n(1-p_n)
\bigl(n^2+2\log_2n+\log_2(1-n^{-2})\bigr)^2\sim n^2.
\]
For $S_n'=(S_n^G,R_n)$ the entropies and information variances add,
by independence. The greedy entropy gap is at most $\log_2e$.
The claimed entropy bound and divergent normalized variance follow.
Each appended alphabet is finite. This does not contradict $W_1$
rigidity: first moments are controlled, but second moments need not be.
\end{proof}

\begin{theorem}[Optimal fixed-confidence information thresholds]\label{thm:thresholds}
For $0<p<1$ define
$Q_n(p)=\inf_S q_p(J_S)$, where the infimum ranges over exact causal
seeds. Then
\[
Q_n(p)=nh+\sqrt n\,q_p(W_{a,b})+o(\sqrt n).
\]
The greedy seed attains this asymptotic for every fixed $p$ simultaneously.
For $b>0$, its explicit coefficient is
\[
q_p(W_{a,b})=
\begin{cases}
b\,\Phi^{-1}\!\left(\dfrac{(a+b)p}{2b}\right),
 &p\le b/(a+b),\\[5pt]
a\,\Phi^{-1}\!\left(1-\dfrac{(a+b)(1-p)}{2a}\right),
 &p\ge b/(a+b).
\end{cases}
\]
For $b=0<a$, it is $a\Phi^{-1}((1+p)/2)$; for $a=b=0$, it is zero.
\end{theorem}
\begin{proof}
Every exact seed satisfies $nh+q_p(Z_n)\le q_p(J_S)$, while the greedy
seed is feasible, so
$nh+q_p(Z_n)\le Q_n(p)\le q_p(J_n^G)$.
At every $p\in(0,1)$ the nondegenerate limiting CDF is continuous and
strictly increasing in a neighborhood of its $p$-quantile. The CDF
limits imply convergence of both normalized bounding quantiles: for
any $\epsilon>0$, the limiting CDF values at $q_p\pm\epsilon$ lie on
opposite sides of $p$, and eventually so do the approximating CDFs.
This proves the assertion by squeezing. Invert \eqref{eq:F} to obtain
the displayed expressions. The deterministic case uses the exact
uniform seed. These are information quantiles, not a claim about exact
fixed-length random-bit generation of arbitrary irrational probabilities.
\end{proof}

\section{Equal positive varentropies: the logarithmic upper bound}\label{sec:fgl}
The main second-order argument above is self-contained. For the sharper
remainder when $a=b=\sigma>0$, we use the following specialization of
Fan--Grama--Liu, Theorem 2.1 \cite{FGL}.

\paragraph{External input (with its exact hypotheses).}
Let $(\xi_i,\mathcal F_i)$ be martingale differences such that, almost
surely,
\[
\sum_{i=1}^n\E(\xi_i^2\mid\mathcal F_{i-1})=1,\qquad
\left|\E(\xi_i^k\mid\mathcal F_{i-1})\right|
\le\tfrac12 k!\epsilon^{k-2}\E(\xi_i^2\mid\mathcal F_{i-1})
\quad(k\ge2),
\]
where $0<\epsilon\le1/2$. There is an absolute constant $C_{\rm FGL}$
such that, for all real $x$,
\begin{equation}\label{eq:fgl}
\left|\Prob\left(\sum_i\xi_i\le x\right)-\Phi(x)\right|
\le C_{\rm FGL}(1+x^2)\epsilon|\ln\epsilon|
\exp\left(-\frac{\widehat x_\epsilon^2}{2}\right),\qquad
\widehat x_\epsilon=\frac{2|x|}{1+\sqrt{1+2\epsilon|x|}}.
\end{equation}
This is that theorem with its quadratic-variation error parameter
$\delta=0$.

\begin{proof}[Proof of \eqref{eq:equalvar}]
Put $B=\max_{A,y}|-\log_2P_A(y)-h|$. For every deterministic policy use
\[
\xi_i=\frac{D_{A_i}(Y_i)}{\sigma\sqrt n},\qquad
\epsilon_n=\frac{B}{\sigma\sqrt n}.
\]
The conditional variance of every $\xi_i$ is exactly $1/n$, regardless
of which action the policy chooses. Also $|\xi_i|\le\epsilon_n$, so
\[
\left|\E(\xi_i^k\mid\mathcal F_{i-1})\right|
\le\epsilon_n^{k-2}\E(\xi_i^2\mid\mathcal F_{i-1})
\le\tfrac12 k!\epsilon_n^{k-2}\E(\xi_i^2\mid\mathcal F_{i-1}).
\]
Thus all the hypotheses hold once $n\ge4B^2/\sigma^2$, with identical
parameters for all policies. Taking the infimum of their CDFs preserves
the same two-sided error about $\Phi$. Consequently
\begin{equation}\label{eq:envfgl}
|G_n(\sigma\sqrt n\,x)-\Phi(x)|
\le C_{\rm FGL}(1+x^2)\epsilon_n|\ln\epsilon_n|
e^{-\widehat x_{\epsilon_n}^2/2}.
\end{equation}
This is a pointwise uniform bound, not merely a separate Wasserstein
estimate for each policy.

The right side has an integrable envelope independent of
$\epsilon_n\le1/2$. Indeed, for $|x|\ge1$,
\[
1+\sqrt{1+2\epsilon_n|x|}
\le1+\sqrt{1+|x|}\le3\sqrt{|x|},\qquad
\widehat x_{\epsilon_n}^2/2\ge2|x|/9.
\]
It follows that
\[
\int_{\R}(1+x^2)e^{-\widehat x_{\epsilon_n}^2/2}\,dx
\le\frac83+2\int_0^\infty(1+x^2)e^{-2x/9}\,dx<400.
\]
Since the standard normal has mean zero, the CDF formula for the
difference of means and \eqref{eq:envfgl} give
\begin{align*}
0\le e_n
&\le\sigma\sqrt n\int_{\R}|G_n(\sigma\sqrt n\,x)-\Phi(x)|\,dx\\
&\le400C_{\rm FGL}B\ln\left(\frac{\sigma\sqrt n}{B}\right).
\end{align*}
The nonnegativity follows also from stochastic dominance over any
single centered policy sum. Lemma~\ref{lem:realization} now gives,
for every integer $n\ge4B^2/\sigma^2$,
\[
0\le C_n^{\causal}-nh
\le400C_{\rm FGL}B\ln\left(\frac{\sigma\sqrt n}{B}\right)+\log_2e.
\]
Absorbing the finitely many smaller horizons into the constant proves
\eqref{eq:equalvar}.
\end{proof}


\section{An explicit error bound for the crash pair}\label{sec:finite}
In this section $P=(1/2,1/4,1/4)$ and $Q=(q,q,1-2q)$.
Let $q$ be the unique root in $(1/3,1/2)$ of
\[
-2q\log_2q-(1-2q)\log_2(1-2q)=3/2.
\]
Uniqueness follows since the derivative is
$2\log_2((1-2q)/q)<0$ on this interval, and the endpoint entropies
are $\log_2 3$ and $1$. Put
\[
a=\tfrac12,\quad b=\sqrt{2q(1-2q)}\log_2\frac q{1-2q},
\qquad c=\sqrt{2/\pi}\,(a-b).
\]
These are exact definitions. The accompanying rational computation gives
\begin{align*}
0.410032428181985576080116712465
&<q<0.410032428181985576080116712466,\\
0.456450775263963373017231101124
&<b<0.456450775263963373017231101125,\\
0.034747254051817864080022137947
&<c<0.034747254051817864080022137948.
\end{align*}

\begin{theorem}[Quantitative crash-pair certificate]\label{thm:finite}
For every integer $n\ge16$,
\begin{equation}\label{eq:finite}
\boxed{\left|\frac{C_n^{\causal}-3n/2}{\sqrt n}-c\right|
\le \frac{17}{2}\frac{\sqrt{\ln n}}{n^{1/4}}
+\frac{3}{\sqrt{n\ln n}}+\frac{2}{\sqrt n}.}
\end{equation}
In particular the unnormalized remainder is
$O(n^{1/4}\sqrt{\ln n})$, with the displayed explicit bound.
\end{theorem}
\begin{proof}
Here is a coarse rational certificate for the constants needed in the
proof, separate from the high-precision decimal coefficient. The positive
logarithm series with the explicit remainder below gives
\[
\frac{1500077}{10^6}<H(.410,.410,.180)<\frac{1500078}{10^6},
\quad
\frac{1497682}{10^6}<H(.411,.411,.178)<\frac{1497683}{10^6}.
\]
Consequently $q_l=.410<q<q_u=.411$. Integer-power comparisons give
\[
\left(\frac{q_l}{1-2q_l}\right)^{16}>2^{19},\qquad
\left(\frac{q_u}{1-2q_u}\right)^4<2^5.
\]
Therefore the defining logarithm for $b$ lies between $19/16$ and $5/4$.
In particular,
\[
\left(\frac9{20}\right)^2
<2q_l(1-2q_u)\left(\frac{19}{16}\right)^2
<b^2
<2q_u(1-2q_l)\left(\frac54\right)^2<\frac14.
\]
For each possible probability $p$ of $P$ or $Q$, the same rational
bracket gives $1/32<p^2<1/2$. Thus
$|{-\log_2p}-3/2|<1$, so $B<1$ in Lemma~\ref{lem:mean}. Since
$|D_A|^3\le B D_A^2$, it follows that $m_3\le1/4$. Also $\pi>2$,
so $K\le20/19$ and
\[
\mathcal C\le\frac{20}{19}
\left(\frac{1/4}{3(9/20)^2}+\frac12\right)
=\frac{4430}{4617}<1.
\]
For a given $n\ge16$ choose
\[
\tau=\sqrt n,\quad \eta_0=n^{-1/4},\quad
R=a\sqrt{2\tau\ln(2/\eta_0)},\quad
s=\sqrt{1+n^{-1/2}},\quad r=R/\sqrt n,\quad T=\sqrt{\ln n}.
\]
The Gaussian Chernoff bound and the prefactors at most two in
\eqref{eq:F} give $\eta\le\eta_0$. Thus, with
$d=\eta+3\mathcal C/\sqrt\tau\le4n^{-1/4}$,
Lemma~\ref{lem:comparison} reads
\[
F((x-r)/s)-d\le G_n(x\sqrt n)\le F((x+r)/s)+d.
\]
Let $Y=Z_n/\sqrt n$ and let $W$ have CDF $F$, so $\E W=c$.
Integrating the last comparison over $[-T,T]$ gives
\begin{equation}\label{eq:core}
\left|\int_{-T}^T\bigl[F(x/s)-G_n(x\sqrt n)\bigr]dx\right|
\le r+2Td.
\end{equation}
To see the constant $r$ explicitly, for any CDF $H$,
$\int_{\R}[H(x)-H(x-r)]dx=r$ when $r\ge0$. Apply this to $H(x)=F(x/s)$
for the upper signed bound and to $H(x+r)-H(x)$ for the lower one.
Each integrand is nonnegative, so restricting its integral to
$[-T,T]$ can only decrease it.

For $t>0$, Lemma~\ref{lem:mean} and \eqref{eq:F} give respectively
\[
\Prob(|Y|>t)\le2e^{-t^2/(2B^2)},\qquad
\Prob(|sW|>t)\le2e^{-t^2/(2a^2s^2)}.
\]
The latter follows by bounding both Gaussian scales by $a$; their
prefactors sum to two. For any $v>0$,
$\int_T^\infty e^{-t^2/(2v)}dt\le(v/T)e^{-T^2/(2v)}$.
Since $B^2\le1$, $a^2s^2\le1$, and
$2B^2+2a^2s^2\le3$, the total contribution outside $[-T,T]$
to the absolute signed difference of the means is at most
\[
\frac3T e^{-T^2/2}=\frac{3}{\sqrt{n\ln n}}.
\]
Consequently \eqref{eq:core} bounds $|\E Y-sc|$ by
$r+2Td+3/\sqrt{n\ln n}$. Next,
$0\le(s-1)c\le1/(2\sqrt n)$, using $c<1$ and
$\sqrt{1+u}-1\le u/2$. For $n\ge16$,
$\ln2\le\tfrac14\ln n$, whence
\[
r=\tfrac12 n^{-1/4}\sqrt{2\ln(2n^{1/4})}
\le\tfrac12 n^{-1/4}\sqrt{\ln n},\qquad
2Td\le8n^{-1/4}\sqrt{\ln n}.
\]
Finally, Lemma~\ref{lem:realization} changes $\E Y$ to
$(C_n^{\causal}-3n/2)/\sqrt n$ by at most
$\log_2e/\sqrt n<3/(2\sqrt n)$; the last inequality follows from
$\ln2>2/3$. Adding these bounds proves \eqref{eq:finite}.
\end{proof}

\begin{center}
\begin{tabular}{@{}rcc@{}}
\toprule
$n$ & Certified lower endpoint & Certified upper endpoint\\
\midrule
$10^{8}$ & 0.000000000000 & 0.399831377112 \\
$10^{12}$ & 0.000000000000 & 0.079430259815 \\
$10^{16}$ & 0.029587976869 & 0.039906531235 \\
$10^{20}$ & 0.034170431871 & 0.035324076233 \\
$10^{24}$ & 0.034684066372 & 0.034810441732 \\
$10^{32}$ & 0.034746524423 & 0.034747983681 \\
$10^{48}$ & 0.034747253962 & 0.034747254142 \\
\bottomrule
\end{tabular}
\end{center}
Every row encloses $(C_n^{\causal}-3n/2)/\sqrt n$. Endpoints are rounded
outward to twelve decimal places. Nonnegativity is also used for the
first two lower endpoints. Already the row $n=10^{16}$ excludes
$c/2=0.0173736270259089\ldots$; later rows narrow around $c$.

\paragraph{Reproduction and arithmetic certificate.}
Run \texttt{python3 certify\_crash\_pair.py}. Only the Python standard
library is required. The output JSON records the coefficient intervals,
all table rows, and the status \texttt{PASS\_EXACT\_RATIONAL\_ARITHMETIC}.
Logarithms are enclosed with 64 terms of the positive series
\[
\ln x=2\sum_{j=0}^{63}\frac{z^{2j+1}}{2j+1}+R,
\qquad z=\frac{x-1}{x+1},\quad
0\le R\le\frac{2z^{129}}{129(1-z^2)}\quad(1\le x\le2),
\]
after powers-of-two range reduction. Entropy signs are checked exactly
at both endpoints at every root-bracketing stage. Square roots use
integer-square-root enclosures; $\pi$ uses Machin's identity with
alternating-series remainders. Every decimal enclosure is rounded
outward. No floating-point value is accepted as an arithmetic premise.
The script certifies this arithmetic, not the universal analytic theorem
by formal verification. The standalone proof above supplies that theorem.


\section{Precise boundaries and remaining questions}\label{sec:open}
\paragraph{All adaptive profiles versus two fixed policies.}
For just the all-$P$ and all-$Q$ policies, their normalized CDF envelope
converges, when $a\ge b>0$, to
\[
F_{\rm fixed}(x)=
\begin{cases}\Phi(x/b),&x\le0,\\\Phi(x/a),&x>0.\end{cases}
\]
Indeed, the ordinary CLT applies to each constant-policy sum, and the
minimum of their two CDFs converges to the displayed minimum. The
uniform Gaussian tails give convergence of means. Its mean is
$(a-b)/\sqrt{2\pi}$, by integrating the two normal tails. The full
adaptive envelope has twice that mean. Thus a ``half constant'' is a
property of this restricted policy family, not of every spectrum method.
A previously quoted bound of order $O(\sqrt n)$ from a ``Theorem 7''
is compatible with identifying a sharp coefficient of $\sqrt n$.
That earlier theorem's complete text is not available here, so this
note does not certify a stronger comparison with it.

\paragraph{The logarithmic regime is not classified.}
For equal positive varentropies we prove an $O(\log(n+1))$ upper bound,
not a matching lower bound or a logarithmic coefficient. A universal
positive lower bound cannot hold: for $P=Q$, the seed consisting of
$n$ independent $P$ samples has entropy $nh$, and one transcript gives
the reverse bound. More generally this is true when the action laws
are relabelings of a common law. Which genuinely distinct pairs have
logarithmic excess remains a question. A difference of third cumulants
is a candidate diagnostic, not a theorem proved in this manuscript.
Uniform control of the adaptive Bellman remainder, rather than an
ordinary independent-sum Edgeworth expansion, is needed to settle it.
No summability of that nonlinear remainder is asserted here.

\paragraph{Constants and degenerate limits.}
The crash-pair finite-horizon estimate has explicit constants. The
$O(\log n)$ application retains the external constant $C_{\rm FGL}$.
The Bellman constant contains $b^{-2}$, so its positive-$b$ bound is
not uniform as $b\downarrow0$; the endpoint is proved separately.
Likewise the many-action theorem assumes a fixed finite family, not
an action family growing with $n$ without uniform hypotheses.
We do not give a sharp unequal-entropy classification in this note.

\paragraph{Computation, horizons, and priorities.}
The algorithm can require exponentially many histories and seed atoms.
There is no polynomial-time claim and no assertion of one
infinite-horizon consistent finite-entropy source. The finite-tree
spectral theorem allows history-dependent kernels, but the explicit
square-root coefficient here is proved for fixed memoryless action
laws of common entropy. Neither a formal proof-assistant certificate
nor independent peer review is asserted. Each literature comparison
concerns the stated models and verified primary-source claims; it is
not a proof of absence of other prior work.

\section{Reproducibility and a small exact optimum}\label{sec:reproduction}
The accompanying standalone Python files use rational arithmetic for
finite-tree reconstruction and certificate checks. For
\[
P=(1/2,\underbrace{1/32,\ldots,1/32}_{16}),\qquad
Q=(\underbrace{1/4,1/4}_{2},\underbrace{1/16,\ldots,1/16}_{8}),
\]
the entropy is $3$ for both actions and the centered information laws
are respectively $\{-2,2\}$ and $\{-1,1\}$, each equiprobable.
At horizon two, the Bellman CDF has masses $1/4$ at centered information
levels $-2,0,1,4$. Its mean is $3/4$ and hence every exact seed has
entropy at least $6+3/4=27/4$. The greedy algorithm constructs a seed
with four atoms of weight $1/16$, sixteen of weight $1/64$, thirty-two
of weight $1/128$, and 256 of weight $1/1024$. Reconstruction of every
history is checked by the independent verifier; its entropy is
$(4+6+7+10)/4=27/4$. Thus this finite instance has exact optimum $27/4$.
The code exports the actual response maps, not just these weights.

The broader verifier enumerates all 128 realized adaptive policies in a
binary depth-three example and compares their laws with the constructed
seed and its Bellman envelope. The additional verifier checks Renyi
moment inequalities, the exact min-entropy endpoint, and the signed
stochastic bounds on three finite examples. These are diagnostic
certificates; the universal proofs are the arguments in this text.
The irrational crash pair is treated by separate rational interval
arithmetic, not by replacing it with a nearby rational law.


\begin{thebibliography}{99}
\bibitem{KDVH17} M. Kocaoglu, A. G. Dimakis, S. Vishwanath, B. Hassibi.
\emph{Entropic Causality and Greedy Minimum Entropy Coupling}.
ISIT 2017, 1465--1469.
\href{https://arxiv.org/abs/1701.08254}{arXiv:1701.08254}.
\bibitem{Li21} C. T. Li.
\emph{Efficient Approximate Minimum Entropy Coupling of Multiple
Probability Distributions}. IEEE Trans. Inform. Theory 67(8),
5259--5268, 2021. \href{https://arxiv.org/abs/2006.07955}{arXiv:2006.07955}.
\bibitem{C22} S. Compton.
\emph{A Tighter Approximation Guarantee for Greedy Minimum Entropy Coupling}.
ISIT 2022. \href{https://arxiv.org/abs/2203.05108}{arXiv:2203.05108}.
\bibitem{CKQGK23} S. Compton, D. Katz, B. Qi, K. Greenewald, M. Kocaoglu.
\emph{Minimum-Entropy Coupling Approximation Guarantees Beyond the
Majorization Barrier}. AISTATS, PMLR 206, 10445--10469, 2023.
\href{https://proceedings.mlr.press/v206/compton23a.html}{Proceedings}.
\bibitem{SY23} Y. Y. Shkel, A. K. Yadav.
\emph{Information Spectrum Converse for Minimum Entropy Couplings and
Functional Representations}. ISIT 2023.
\href{https://arxiv.org/abs/2305.05745}{arXiv:2305.05745}.
\bibitem{MWW25} Y.-J. Ma, F. Wang, X.-Y. Wu.
\emph{Efficient approximate Minimum-Renyi entropy couplings}.
Discrete Contin. Dyn. Syst. S 18(10), 2924--2936, 2025.
\href{https://doi.org/10.3934/dcdss.2025098}{doi:10.3934/dcdss.2025098}.
\bibitem{MWWC25} Y.-J. Ma, F. Wang, X.-Y. Wu, K.-Y. Cai.
\emph{An Explicit Description of Extreme Points of the Set of Couplings
with Given Marginals: with Application to Minimum-Entropy Coupling Problems}.
2025. \href{https://arxiv.org/abs/2505.12227}{arXiv:2505.12227}.
\bibitem{Peng07} S. Peng.
\emph{Law of Large Numbers and Central Limit Theorem under Nonlinear
Expectations}. 2007.
\href{https://arxiv.org/abs/math/0702358}{arXiv:math/0702358}.
\bibitem{FPSS18} X. Fang, S. Peng, Q.-M. Shao, Y. Song.
\emph{Limit theorems with rate of convergence under sublinear expectations}.
Preprint, version 2, 2018.
\href{https://arxiv.org/abs/1711.10649}{arXiv:1711.10649}.
\bibitem{PZ19} S. Peng, Q. Zhou.
\emph{A hypothesis-testing perspective on the G-normal distribution theory}.
Preprint, 2019.
Theorem 2. \href{https://arxiv.org/abs/1909.03530}{arXiv:1909.03530}.
\bibitem{FGL} X. Fan, I. Grama, Q. Liu.
\emph{Nonuniform Berry-Esseen bounds for martingales with applications
to statistical estimation}. Theorem 2.1, conditions (A1)--(A2).
\href{https://arxiv.org/abs/1608.05217}{arXiv:1608.05217}.
\end{thebibliography}
\end{document}
